\documentclass[a4paper]{arrowhead}

\usepackage[yyyymmdd]{datetime}
\usepackage{etoolbox}
\usepackage[utf8]{inputenc}
\usepackage[T1]{fontenc}
\usepackage{multirow}
\usepackage{longtable}
\usepackage{tikz}
\usepackage{adjustbox}
\AtBeginEnvironment{longtable}{\small}
\BeforeBeginEnvironment{tikzpicture}{\begin{adjustbox}{max width=\textwidth}}
\AfterEndEnvironment{tikzpicture}{\end{adjustbox}}
\usetikzlibrary{positioning,shapes.multipart,arrows.meta}

\renewcommand{\dateseparator}{-}
\setlength{\emergencystretch}{3em}
\lstset{language=HTTP,breakatwhitespace=false,columns=fullflexible}

%% Special references
\newcommand{\mref}[2][]{{\textcolor{ArrowheadPurple}{\hyperref[sec:model:#2]{\ifstrempty{#1}{#2}{#1}}}}}
\newcommand{\pref}[1]{{\textcolor{ArrowheadGrey}{\hyperref[sec:model:primitives:#1]{#1}}}}
%%

\begin{document}
\sloppy

%% Arrowhead Document Properties
\ArrowheadTitle{Authentication}
\ArrowheadType{System Design Description}
\ArrowheadTypeShort{SysDD}
\ArrowheadVersion{5.2.2}
\ArrowheadDate{2026-11-07}
\ArrowheadAuthor{Jerker Delsing}
\ArrowheadStatus{RELEASE CANDIDATE}
\ArrowheadContact{jerker.delsing@ltu.se}
\ArrowheadFooter{\href{www.arrowhead.eu}{www.arrowhead.eu}}
\ArrowheadSetup
%%

%% Front Page
\begin{center}
  \vspace*{1cm}
  \huge{\arrowtitle}

  \vspace*{0.2cm}
  \LARGE{\arrowtype}
  \vspace*{1cm}

  \vspace*{\fill}

  % Front Page Image

  \vspace*{1cm}
  \vspace*{\fill}

  % Front Page Abstract
  \begin{abstract}
    This document describes the design and implementation of the Authentication core system of Arrowhead 5.2.2 in Java Spring (\texttt{ah5-core-java-spring/authentication}).
  \end{abstract}

  \vspace*{1cm}
\end{center}

\newpage
%%

%% Table of Contents
\tableofcontents
\newpage
%%

\section{Overview}
\label{sec:overview}

This document describes the Authentication system implementation, which realizes the Authentication SysD \cite{auth-sysd}. Authentication is the identity provider of an Arrowhead 5 local cloud. It stores the identities of the systems with their credentials and sysop flag, lets systems log in and out, and issues and verifies the identity tokens with which systems authenticate towards each other.
In Section \ref{Implementation}, we describe implementation details of the system.

\newpage

\section{Implementation}
\label{Implementation}

This implementation is based on the SysD document Authentication \cite{auth-sysd}, the upstream Eclipse Arrowhead Java Spring implementation 5.2.2 \cite{ah5java} and the local changes listed in Section \ref{sec:local}.

\subsection{Implementation language and tools}
\begin{itemize}
\item Java 21, Spring Boot 3.4.9 (web, security, data JPA, Quartz), springdoc-openapi 2.8 (Swagger UI), Log4j 2.
\item Maven module \texttt{authentication} of \texttt{ah5-core-java-spring}; it depends on \texttt{arrowhead-common-utils} and \texttt{arrowhead-data-transfer-objects} 5.2.2 from the local \texttt{ah5-common-java-spring} (Eclipse Paho MQTT v3 1.2.5, RabbitMQ AMQP client 5.22.0).
\item Database: MariaDB/MySQL (MariaDB Connector/J; MySQL 8.4 server in the test cloud), schema validated by Hibernate (\texttt{ddl-auto=validate}); SQL scripts in \texttt{src/main/resources/database}.
\item Stateless request processing; all state is in the database. Scheduled jobs use Quartz.
\end{itemize}

\subsection{Functional properties implementation}

The Java packages below \texttt{eu.arrowhead} implement the system:

\noindent\begin{longtable}{| p{5.6cm} | p{9.6cm} |} \hline
\rowcolor{gray!33} Package & Responsibility \\ \hline \endhead
\texttt{api.\allowbreak{}http,\allowbreak{} api.\allowbreak{}mqtt} & Controllers and topic handlers of identity, identity\allowbreak{}Management and generalManagement. \\ \hline
\texttt{method,\allowbreak{} method.\allowbreak{}password} & Pluggable authentication methods; the password method with its validator, normalizer and database service. \\ \hline
\texttt{service} & Identity and management logic; validation and normalization. \\ \hline
\texttt{http.\allowbreak{}filter,\allowbreak{} mqtt.\allowbreak{}filter} & \texttt{Internal\allowbreak{}Authentication\allowbreak{}Filter}: token check in the own database, token-to-certificate binding (local). \\ \hline
\texttt{jpa} & System, password authentication, active session and log entities. \\ \hline
\texttt{quartz} & Cleaner job for expired sessions. \\ \hline
\end{longtable}

Every request passes the common filter chain: authentication filter (by \texttt{authentication.policy}), optional blacklist filter, management service filter (for management services), then the controller, validation and normalization, the service logic and the database service. The same logic is reachable through the MQTT topic handlers.

\subsubsection{Data handled and stored}

Database \texttt{ah\_authentication}:

\noindent\begin{longtable}{| p{6.4cm} | p{8.8cm} |} \hline
\rowcolor{gray!33} Tables & Content \\ \hline \endhead
\texttt{system\_\allowbreak{}} & Identities with sysop flag and authentication method. \\ \hline
\texttt{password\_\allowbreak{}authentication} & Password hashes. \\ \hline
\texttt{active\_\allowbreak{}session} & Sessions: token, login and expiry time. \\ \hline
\texttt{logs} & Log entries for \texttt{get-\allowbreak{}log}. \\ \hline
\end{longtable}

The data models of the interfaces are the Java records of \texttt{eu.arrowhead.dto} (common library); the IDD documents name them without the \texttt{DTO} suffix. The records used by this system are: \texttt{Identity\allowbreak{}Change\allowbreak{}Request\allowbreak{}DTO}, \texttt{Identity\allowbreak{}List\allowbreak{}Mgmt\allowbreak{}Create\allowbreak{}Request\allowbreak{}DTO}, \texttt{Identity\allowbreak{}List\allowbreak{}Mgmt\allowbreak{}Response\allowbreak{}DTO}, \texttt{Identity\allowbreak{}List\allowbreak{}Mgmt\allowbreak{}Update\allowbreak{}Request\allowbreak{}DTO}, \texttt{Identity\allowbreak{}Login\allowbreak{}Response\allowbreak{}DTO}, \texttt{Identity\allowbreak{}Mgmt\allowbreak{}Request\allowbreak{}DTO}, \texttt{Identity\allowbreak{}Mgmt\allowbreak{}Response\allowbreak{}DTO}, \texttt{Identity\allowbreak{}Query\allowbreak{}Request\allowbreak{}DTO}, \texttt{IdentityRequestDTO}, \texttt{Identity\allowbreak{}Session\allowbreak{}List\allowbreak{}Mgmt\allowbreak{}Response\allowbreak{}DTO}, \texttt{Identity\allowbreak{}Session\allowbreak{}Query\allowbreak{}Request\allowbreak{}DTO}, \texttt{Identity\allowbreak{}Session\allowbreak{}Response\allowbreak{}DTO}, \texttt{Identity\allowbreak{}Verify\allowbreak{}Response\allowbreak{}DTO}, \texttt{KeyValuesDTO}, \texttt{LogEntryDTO}, \texttt{Log\allowbreak{}Entry\allowbreak{}List\allowbreak{}Response\allowbreak{}DTO}, \texttt{LogRequestDTO}, \texttt{PageDTO}.

\subsection{Non functional properties implementation}

\subsubsection{Security}

TLS and client authentication are configured with the Spring \texttt{server.ssl.*} properties (PKCS\#12 keystore and truststore). The authentication filters of the common library implement the four authentication policies; \texttt{TokenCertificateBinding} binds identity tokens to the client certificate and \texttt{CrlRevocationConfig} enables CRL checking in Tomcat (both local). The private key of the keystore is also used to sign the system's outgoing requests where required.

\subsubsection{Power management}

None; the system is a server process without power states. The \texttt{sysMonitor} \texttt{power} observation always reports line supply.

\subsubsection{Internal monitoring}

Log4j 2 logs to stdout and rolling files; log entries are also written to the \texttt{logs} table and served by \texttt{generalManagement.get-log}. The native \texttt{sysMonitor} publisher (\texttt{eu.arrowhead.common.sysmonitor}) publishes \texttt{echo}, \texttt{power} and \texttt{status} every 30 s.

\subsubsection{Configuration}

Configuration is read from \texttt{application.properties} and environment variables (Spring relaxed binding, e.g. \texttt{SERVICE\_DISCOVERY\_POLICY}); secrets come from files through \texttt{SPRING\_CONFIG\_IMPORT=configtree:/run/secrets/}. The properties are listed in the SysD. \texttt{generalManagement.get-config} returns the public ones.

\subsection{Deployment}

The test cloud (\texttt{Arrowhead-test-cloud-java}) builds the image \texttt{arrowhead-\allowbreak{}authentication:5.\allowbreak{}2.\allowbreak{}2-\allowbreak{}sysmonitor} from the released image \texttt{aitiaiiot/\allowbreak{}arrowhead-\allowbreak{}authentication:5.\allowbreak{}2.\allowbreak{}2} with the locally built \texttt{arrowhead-\allowbreak{}authentication-\allowbreak{}5.\allowbreak{}2.\allowbreak{}2.\allowbreak{}jar} (\texttt{java-patch/Dockerfile}). The container runs read-only with the configuration baked into \texttt{/app/config}, only its keystores mounted, and publishes port 8444 on the loopback address of the host. Its database runs in a separate container.

\subsection{Local modifications}
\label{sec:local}

Compared to upstream 5.2.2, this workspace changes:
\begin{itemize}
  \item Startup is refused when the effective configuration is below the secure profile (\texttt{Security\allowbreak{}Profile\allowbreak{}Guard}).
  \item Login, logout and credential change accepted only with the client certificate of the same system; token-to-certificate binding in the authentication filter.
  \item \texttt{http.\allowbreak{}collector.\allowbreak{}bootstrap.\allowbreak{}services} includes \texttt{authorization}: the management filter finds Consumer\allowbreak{}Authorization in the ServiceRegistry, not through orchestration.
  \item Seeded identities: ServiceRegistry and Sysop passwords are changed to per-cloud secrets after the first start (\texttt{provision-\allowbreak{}authentication-\allowbreak{}identities.\allowbreak{}sh}); only Sysop is a sysop.
  \item Native \texttt{sysMonitor} publisher, CRL checking, MariaDB driver (common library and image patch).
\end{itemize}

\subsection{Verification}

The module has 17 JUnit test classes; the regression report set recorded on 2026-10-20 covers 374 tests without failures. Live-cloud smoke tests are a separate verification step; the plant-management migration is not claimed as deployed by this update.

\newpage

\section{Services}
\label{sec:services}

The implementation of the services is based on the SD and IDD documents referenced below.

\noindent\begin{longtable}{| p{5.4cm} | p{2.2cm} | p{2.2cm} | p{4.8cm} |} \hline
\rowcolor{gray!33}\textbf{Services produced} & \textbf{SysD ref} & \textbf{SD ref} & \textbf{IDD ref} \\ \hline \endhead
\texttt{identity} & \cite{auth-sysd} & \cite{auth-identity-sd} & \cite{auth-identity-http-idd}, \cite{auth-identity-mqtt-idd} \\ \hline
\texttt{identityManagement} & \cite{auth-sysd} & \cite{auth-identityManagement-sd} & \cite{auth-identityManagement-http-idd}, \cite{auth-identityManagement-mqtt-idd} \\ \hline
\texttt{generalManagement} & \cite{auth-sysd} & \cite{auth-generalManagement-sd} & \cite{auth-generalManagement-http-idd}, \cite{auth-generalManagement-mqtt-idd} \\ \hline
\texttt{sysMonitor} & \cite{auth-sysd} & \cite{auth-sysMonitor-sd} & \cite{auth-sysMonitor-amqp-idd} \\ \hline
&&& \\ \hline
\rowcolor{gray!33}\textbf{Services consumed} & \textbf{SysD ref} & \textbf{SD ref} & \textbf{IDD ref} \\ \hline
serviceDiscovery & \cite{sr-sysd} & \cite{sr-serviceDiscovery-sd} & -- \\ \hline
authorization & \cite{ca-sysd} & \cite{ca-authorization-sd} & -- \\ \hline
sysMonitor (publisher) & AMQP broker (SysMonitor) & -- & -- \\ \hline
\end{longtable}


\bibliographystyle{IEEEtran}
\bibliography{bibliography}

\newpage

\section{Revision History}
\subsection{Amendments}

\noindent\begin{tabularx}{\textwidth}{| p{1cm} | p{3cm} | p{2cm} | X | p{4cm} |} \hline
\rowcolor{gray!33} No. & Date & Version & Subject of Amendments & Author \\ \hline

1 & 2026-11-07 & \arrowversion & Document created from the Arrowhead LaTeX templates for the 5.2.2 Java Spring implementation, including the local extensions of this workspace. & Jerker Delsing \\ \hline
\end{tabularx}

\subsection{Quality Assurance}

\noindent\begin{tabularx}{\textwidth}{| p{1cm} | p{3cm} | p{2cm} | X |} \hline
\rowcolor{gray!33} No. & Date & Version & Approved by \\ \hline

1 & & \arrowversion  &  \\ \hline

\end{tabularx}

\end{document}
